%======================================================================
\section{Preliminaries}\label{sec:prelim}

Throughout, $\Sigma$ denotes a finite alphabet, words are finite strings over $\Sigma$ (including the empty word), and
$|w|$ is the length of a word $w$. For a ring $R$ we write $[w]$ for the image of a word in an algebra presented on
$\Sigma$.

\subsection{Word presentations over the integers}

An \emph{elementary replacement} is a relation $u=v$ between nonempty words of the same length; it may be used in either
direction inside a word. A \emph{zero declaration} is a relation $w=0$ with $w$ nonempty. For a commutative ring
$\Lambda$ and a finite set $\mathcal S$ of elementary replacements and zero declarations, let $A(\Lambda)$ be the unital
associative $\Lambda$-algebra generated by $\Sigma$ subject to $\mathcal S$. Thus
$A(\Lambda)=\Lambda\langle\Sigma\rangle/\mathcal J_\Lambda$, where $\mathcal J_\Lambda$ is the two-sided ideal generated
by the binomials $u-v$ and the words $w$ listed in $\mathcal S$.

The \emph{replacement graph} has the words as vertices and an edge between $xuy$ and $xvy$ for every listed replacement
$u=v$ and all words $x,y$. A connected component of this graph is a \emph{replacement class}. A class is \emph{zero} if
one of its words contains a zero-declared word as a substring, and \emph{nonzero} otherwise. The following lemma is
\cite[Lemma~2.1]{OAI247periodic}, stated there over a field; we need it over $\Z$, so that a single presentation can be
used in every characteristic.

\begin{lemma}\label{lem:word-basis}
The images of the nonzero replacement classes form a basis of $A(\Z)$ as a free $\Z$-module, and the zero classes map to
$0$. For every commutative ring $\Lambda$, the natural map $\Lambda\otimes_{\Z}A(\Z)\to A(\Lambda)$ is an isomorphism, so
the images of the nonzero classes form a $\Lambda$-basis of $A(\Lambda)$. In particular, whether a word is zero does not
depend on $\Lambda$, and $A(\Lambda)$ is graded by word length with $A(\Lambda)_0=\Lambda$.
\end{lemma}

\begin{proof}
Let $V$ be the free $\Z$-module on all words, and let $V_1\subseteq V$ be the span of the differences $x-y$ for edges
$\{x,y\}$ of the replacement graph. A finite combination $\sum_w c_ww$ lies in $V_1$ if and only if, for every
replacement class $C$, the coefficient sum $\sum_{w\in C}c_w$ vanishes. Indeed, edge differences have this property. Conversely, within one class each word is
joined to a fixed base word by a path, and the difference of its endpoints telescopes into edge differences. Hence $V/V_1$
is free on the classes. The two-sided ideal of $\Z\langle\Sigma\rangle$ generated by the listed binomials is $V_1$, since
$V_1$ is spanned by the products $x(u-v)y$.

Multiplication of words on either side maps replacement classes to replacement classes, and maps zero classes to zero
classes. Let $V_0\subseteq V/V_1$ be the span of the zero classes. It is a two-sided ideal, it contains the images of the
zero-declared words, and it is contained in the ideal they generate, because every word of a zero class is equivalent,
modulo $V_1$, to a word $xwy$ with $w$ zero-declared. Hence $A(\Z)=(V/V_1)/V_0$ is free on the nonzero classes.

For an arbitrary $\Lambda$, right exactness of $\Lambda\otimes_{\Z}-$ applied to the presentation of $A(\Z)$ gives
$A(\Lambda)\cong\Lambda\otimes_{\Z}A(\Z)$, and the basis is carried along. Every relation is homogeneous of positive
length, so word length induces the grading, and the empty word spans the degree-zero part.
\end{proof}

Let $A_+(\Lambda)\subseteq A(\Lambda)$ be the span of the nonempty words. The proof of \cite[Corollary~3.8]{OAI247periodic}
shows, over any commutative ring, that the unital ideal $\mathcal J_\Lambda$ coincides with the two-sided ideal of the free
nonunital algebra $\Lambda\langle\Sigma\rangle_+$ generated by the same finite set of relations. Both ideals are spanned
by the products $xsy$, with $s$ a defining relation and $x,y$ possibly empty words. Consequently:

\begin{lemma}\label{lem:nonunital}
$A_+(\Lambda)$ is finitely presented as a nonunital $\Lambda$-algebra, by the generators $\Sigma$ and the relations
$\mathcal S$, and $A(\Lambda)=\Lambda\oplus A_+(\Lambda)$ is its unitization.
\end{lemma}

\subsection{Steinberg groups}

For a unital ring $R$ and $n\geq3$, the Steinberg group $\St_n(R)$ is generated by symbols $x_{ij}(a)$, where
$1\leq i\neq j\leq n$ and $a\in R$, subject to
\begin{align}
 x_{ij}(a+b)&=x_{ij}(a)x_{ij}(b),\label{eq:st-add}\\
 [x_{ij}(a),x_{kl}(b)]&=1\qquad (j\neq k,\ i\neq l),\label{eq:st-comm}\\
 [x_{ij}(a),x_{jk}(b)]&=x_{ik}(ab)\qquad (i,j,k\text{ pairwise distinct}),\label{eq:st-prod}
\end{align}
where $[g,h]=ghg^{-1}h^{-1}$. The assignment $x_{ij}(a)\mapsto e_{ij}(a)$, the elementary matrix with $a$ in position
$(i,j)$, defines a surjection $\pi_n:\St_n(R)\to\E_n(R)$ onto the elementary subgroup of $\GL_n(R)$. Its kernel is
$J_n(R)$. Write $\St(R)$, $\E(R)$ and $J(R)$ for the direct limits over $n$.

\subsection{The finite-rule hierarchy}

We use OpenAI's hierarchy of finite rule tables~\cite[Section~2]{OAI247periodic} without change. We recall what it
provides. Fix the absolute constants $c=10000$ and $k=10$ of \cite[Proposition~2.2]{OAI247periodic}, an integer
\begin{equation}\label{eq:D}
 D>10(c+1)(k+1),
\end{equation}
and an integer $h_0$, discussed in Remark~\ref{rem:parameters}. For $j\geq0$ put
\[
 h_j=h_0D^j,\qquad q_j=2^{h_j},\qquad L_j=2q_j+1,
\]
and let $\Sigma_j$ and $\mathcal S_j$ be the alphabet and rule table at parameter $h_j$. Write $A_j(\Lambda)$ for the
corresponding algebra.

\begin{proposition}[{\cite[Proposition~2.2]{OAI247periodic}}]\label{prop:hierarchy}
For every $j\geq0$ the following hold.
\begin{enumerate}[label=\textup{(H\arabic*)}]
\item\label{H:encoding} $|\Sigma_j|\leq2^{ch_j}$, and every rule of $\mathcal S_j$ has length at most $k$. There is a
  substitution $E_j:\Sigma_{j+1}\to\Sigma_j^{L_j}$ that carries every rule of $\mathcal S_{j+1}$ to a consequence of
  $\mathcal S_j$. It therefore induces a unital homomorphism $E_j:A_{j+1}(\Lambda)\to A_j(\Lambda)$. If $w$ is a nonzero
  word at level $j+1$, then $E_j(w)$ is a nonzero word at level $j$. Every level has a nonzero one-letter word.
\item\label{H:tokens} A designated set of letters marks \emph{cuts} immediately before them. There are a bound $B_j$ and a
  finite list $\mathcal T_j$ of \emph{tokens} such that every sufficiently long nonzero word at level $j$ has a unique
  decomposition $p\,t_1\cdots t_r\,s$ with $t_i\in\mathcal T_j$ and $|p|,|s|\leq B_j$. Here $p$ precedes the first cut, $s$
  starts at the last cut, and the $t_i$ are the intervals between consecutive cuts. For each $t\in\mathcal T_j$ there is a
  word $u_t$ over $\Sigma_{j+1}$ with $1\leq|u_t|\leq k$, $|t|=L_j|u_t|$ and $[t]=E_j([u_t])$.
\item\label{H:runs} Some letters are \emph{palette letters}, with a color in a set of $q_j$ colors. An \emph{eligible run}
  is a maximal string of consecutive palette letters with a common tuple of non-color fields and no control head. Every
  adjacent transposition inside an eligible run is an elementary replacement. Every eligible run in a nonzero word has
  pairwise distinct colors. Every sufficiently long nonzero word contains an eligible run of length at least $q_j/100$.
  Eligibility and the endpoints of maximal eligible runs depend only on the non-color fields.
\end{enumerate}
\end{proposition}

\begin{remark}[Independence of the field]\label{rem:char-free}
OpenAI state explicitly that a field is not part of the rule table: the same word relations are used over every
field~\cite[Section~2]{OAI247periodic}. Properties \ref{H:encoding}--\ref{H:runs} concern words, the replacement graph
and nonzero classes. By Lemma~\ref{lem:word-basis} these do not depend on the coefficient ring. We have checked that the
proofs in \cite[Section~2]{OAI247periodic} use no property of $\F_2$. Thus Proposition~\ref{prop:hierarchy} holds verbatim
over $\Z$ and over every field.
\end{remark}

\begin{remark}[Freedom in the parameters]\label{rem:parameters}
In \cite[Section~2.5]{OAI247periodic} the integer $D$ is fixed before the program fixed point. Afterwards $h_0$ is chosen
so large that, for every $h\geq h_0$, all field and layout bounds hold and every relevant computation finishes before the
tableau clock $T=2^{\lfloor h/5\rfloor}$ expires. Hence the conclusions of Proposition~\ref{prop:hierarchy} hold for every
integer $h_0$ exceeding a threshold that depends only on $D$ and the fixed program. Likewise $D$ may be any integer
satisfying~\eqref{eq:D}. Neither Theorem~A nor Theorem~B needs this freedom: we fix $D$ and $h_0$ once and for all, and
Section~\ref{sec:detection} handles every prime with the same table.
\end{remark}

%======================================================================
\section{The algebra in every positive characteristic}\label{sec:algebra}

We follow \cite[Section~3]{OAI247periodic}. Only one step depends on the characteristic: the cancellation of a symmetric
matrix sum inside a palette run. We replace it by a bound that depends on $p$ (Lemma~\ref{lem:symmetric}), and then check that the resulting
factor $p-1$ is absorbed by the hierarchy at a finite level.

Throughout this section $K$ is a field of characteristic $p>0$. The algebras $A_j(K)=K\otimes_{\Z}A_j(\Z)$ are those of
Section~\ref{sec:prelim}, and $E_j:A_{j+1}(K)\to A_j(K)$ is the homomorphism of~\ref{H:encoding}. We identify
$M_d(K)\otimes_KA_j(K)=M_d(A_j(K))$. A \emph{letter test} and a \emph{chunk test} of dimension $d$ (and label bound $b$)
at level $j$ are defined exactly as in \cite[Definition~3.1]{OAI247periodic}, with matrix weights in $M_d(K)$. Their
length-$n$ sums $S_{j,n}(M)$ and $C_{j,n}(\mathcal C,W)$ are as defined there. A test
\emph{eventually vanishes} if its length-$n$ sum is zero for all sufficiently large $n$.

\subsection{Symmetric sums in characteristic \texorpdfstring{$p$}{p}}

\begin{lemma}\label{lem:symmetric}
Let $X_1,\dots,X_m\in M_d(K)$, where $K$ has characteristic $p$, and put
\[
 \mathcal A_m(X_1,\dots,X_m)=\sum_{\sigma\in S_m}X_{\sigma(1)}\cdots X_{\sigma(m)}.
\]
If $m>(p-1)d^2$, then $\mathcal A_m(X_1,\dots,X_m)=0$.
\end{lemma}

\begin{proof}
Choose a basis $Y_1,\dots,Y_{d^2}$ of $M_d(K)$ and write $X_t=\sum_ic_{t,i}Y_i$. Since $\mathcal A_m$ is multilinear,
\[
 \mathcal A_m(X_1,\dots,X_m)=\sum_{(i_1,\dots,i_m)}\Bigl(\prod_{t=1}^mc_{t,i_t}\Bigr)
 \sum_{\sigma\in S_m}Y_{i_{\sigma(1)}}\cdots Y_{i_{\sigma(m)}} .
\]
Fix an index tuple $(i_1,\dots,i_m)$ and let $n_1,\dots,n_{d^2}$ be the multiplicities with which the basis indices
occur in it, so $\sum_in_i=m$. As $\sigma$ ranges over $S_m$, each distinct arrangement of these indices occurs exactly
$\prod_in_i!$ times. The inner sum therefore equals $\prod_in_i!$ times a sum of words in the $Y_i$. Since $m>(p-1)d^2$,
some $n_i\geq p$, so $p$ divides $\prod_in_i!$ and the inner sum vanishes in characteristic $p$.
\end{proof}

For $p=2$ this is the bound $m>d^2$ of \cite[Lemma~3.2]{OAI247periodic}. For $d\geq2$ the bound is far from sharp.
Zalesskii proved that $\mathcal A_{dp}$ vanishes identically on $M_d(K)$ in characteristic $p$, and that $\mathcal A_k$
does not for $k<dp$~\cite{Zalesskii1985}; we take the statement from~\cite{Valenti1998}. Grouping the terms of
$\mathcal A_{m+1}$ by their last factor writes it as a sum of terms $\mathcal A_m(\dots)X_t$, so $\mathcal A_m$ then
vanishes for every $m\geq dp$. Any bound that is finite for each $p$ and $d$ suffices below, and we keep
Lemma~\ref{lem:symmetric} because its proof is two lines. Its bound is sharp for $d=1$: for scalars equal
to $1$ the sum is $m!$, which is nonzero modulo $p$ when $m\leq p-1$. No signs or divisions occur. The quantity that
vanishes is a positive integer multiplicity, and this mechanism has no analogue in characteristic zero
(Section~\ref{sec:charzero}).

\subsection{Cancellation and its propagation}

\begin{lemma}[Direct cancellation in characteristic $p$]\label{lem:direct}
Fix $j,d,b$. If
\begin{equation}\label{eq:direct-bound}
 \frac{q_j}{100}>b\bigl((p-1)d^2+3\bigr),
\end{equation}
then every chunk test of dimension $d$ and label bound $b$ at level $j$ eventually vanishes over $K$.
\end{lemma}

\begin{proof}
The proof of \cite[Lemma~3.3]{OAI247periodic} applies with one change. It partitions the nonzero summands of a long
length-$n$ sum into free orbits of the symmetric group $S_m$. The group acts by permuting the $m$ complete chunks (the
\emph{packet}) inside the first eligible run of length at least $q_j/100$. Each orbit contributes
$U\,\mathcal A_m(X_1,\dots,X_m)\,V\otimes a$. That partition uses only~\ref{H:runs} and the replacement graph, and no
property of the field.

If the chosen run has length $r$ and its packet has $m$ chunks, then $r\leq mb+2b$, since at most two chunks cross its
endpoints. By~\eqref{eq:direct-bound}, $m\geq r/b-2>(p-1)d^2+1$. Hence every orbit contribution vanishes by
Lemma~\ref{lem:symmetric}, in place of \cite[Lemma~3.2]{OAI247periodic}.
\end{proof}

The next two results are \cite[Lemmas~3.4 and~3.5]{OAI247periodic}. Their proofs are linear algebra together
with~\ref{H:tokens}, and are valid over every field.

\begin{lemma}[{Reduction by the cuts~\cite[Lemma~3.4]{OAI247periodic}}]\label{lem:cut}
If every chunk test of dimension $d$ and label bound $k$ at level $j+1$ eventually vanishes, then every letter test of
dimension $d$ at level $j$ eventually vanishes.
\end{lemma}

\begin{lemma}[{A matrix automaton for chunks~\cite[Lemma~3.5]{OAI247periodic}}]\label{lem:automaton}
Let a chunk test over an alphabet of size $s$ have dimension $d$ and label bound $b$. There is a letter test of dimension
at most $b\,s^bd$ whose length-$n$ sum contains the chunk sum as a fixed $d\times d$ block, for every $n$.
\end{lemma}

\begin{proposition}[Eventual vanishing of every fixed test]\label{prop:all-tests}
Over every field $K$ of characteristic $p>0$, every letter test at level $0$ eventually vanishes, and so does every chunk
test at level $0$ with any fixed label bound.
\end{proposition}

\begin{proof}
We follow \cite[Proposition~3.6]{OAI247periodic}. Fix a starting dimension $d$, and define
$d_j=d\prod_{i=1}^{j-1}(k|\Sigma_i|^k)$. Lemmas~\ref{lem:cut} and~\ref{lem:automaton} reduce eventual vanishing of a
dimension-$d$ letter test at level $0$ to eventual vanishing of all dimension-$d_J$ chunk tests of label bound $k$ at
level $J$, for any $J\geq1$. As in the proof of \cite[Proposition~3.6]{OAI247periodic},
\[
 \log_2(d_j^2)\leq2\log_2d+2(j-1)\log_2k+\frac{2kc}{D-1}(h_j-h_1),
\]
and $\delta=1-2kc/(D-1)>0$ by~\eqref{eq:D}. Hence $h_j-\log_2(d_j^2)\to\infty$, and there is a finite $J$ with
\[
 2^{h_J}>400k(p-1)d_J^2\geq100k\bigl((p-1)d_J^2+3\bigr),
\]
where the second inequality uses $(p-1)d_J^2\geq1$. This is~\eqref{eq:direct-bound} with $b=k$. Lemma~\ref{lem:direct}
therefore applies at level $J$, and the finitely many reductions pull the vanishing back to level $0$, exactly as in
\cite{OAI247periodic}. The only difference from characteristic two is the factor $p-1$, which changes the finite stopping
level $J$ but not the presentation. A level-zero chunk test of any label bound is a block of a letter test, by
Lemma~\ref{lem:automaton}.
\end{proof}

\subsection{Matrix nilness and the proof of Theorem~A}

Let $\Sigma=\Sigma_0$ and $\mathcal S=\mathcal S_0$ be the level-zero alphabet and rule table, with $D$ and $h_0$ fixed
once and for all. This is the presentation of Theorem~A; it is written over $\Z$. For a field $K$ put $R_K=A_0(K)$ and
let $I_K=A_+(K)$ be its positive-degree ideal.

\begin{theorem}\label{thm:matrix-nil}
Let $K$ be a field of positive characteristic. Then $M_\ell(I_K)$ is nil for every $\ell\geq1$: every matrix with
entries in $I_K$ is nilpotent, with an index depending on the matrix.
\end{theorem}

\begin{proof}
This is the coefficient extraction of \cite[Theorem~3.7]{OAI247periodic}, run over an arbitrary field of characteristic
$p$. Write $T\in M_\ell(I_K)$ as $T=\sum_{\gamma}A_\gamma\otimes[v_\gamma]$, with $A_\gamma\in M_\ell(K)$ and
$|v_\gamma|\geq1$, and let $b=\max_\gamma|v_\gamma|$. Let $z$ be an indeterminate.

Apply Proposition~\ref{prop:all-tests} over the field $K(z)$, which has characteristic $p$, to the chunk test with labels
$v_\gamma$ and weights $zA_\gamma$. Its length-$n$ sum is the image of a polynomial $P_n(z)\in M_\ell(R_K)[z]$, as in the proof of
\cite[Theorem~3.7]{OAI247periodic}, and it vanishes for $n\geq N$. The map
$M_\ell(R_K)[z]\to M_\ell(K(z)\otimes_KR_K)$ is injective: expand coefficients in a $K$-basis of $R_K$ and use the
injectivity of $K[z]\to K(z)$. So $P_n(z)=0$ for $n\geq N$.

For $a\geq1$ we have $T^a=\sum_{n=a}^{ab}[z^a]P_n(z)$, since each factor has a label of length between $1$ and $b$.
Hence $T^a=0$ for $a\geq N$.
\end{proof}

The proof establishes cancellation afresh over every field, including $K(z)$. It does not rely on the general
principle, false for arbitrary nil algebras, that nilness survives extension of scalars.

\begin{proof}[Proof of Theorem~A]
\begin{itemize}[leftmargin=*]
\item \emph{Finite presentation.} By Lemma~\ref{lem:nonunital}, $I_K$ is finitely presented as a nonunital algebra by
  $\Sigma$ and $\mathcal S$, and $R_K=K\oplus I_K$ is its unitization, finitely presented by the same data.
\item \emph{Infinite dimension.} By~\ref{H:encoding}, the words $w_j=E_0E_1\cdots E_{j-1}(a_j)$, with $a_j$ a nonzero
  letter at level $j$, are nonzero and have strictly increasing lengths $L_0\cdots L_{j-1}$. By
  Lemma~\ref{lem:word-basis} they are nonzero over every field and linearly independent, so $I_K$ is
  infinite-dimensional.
\item \emph{Not nilpotent.} If $I_K^m=0$, every word of length at least $m$ would vanish, and $I_K$ would be spanned by the
  finitely many words of length less than $m$, a contradiction.
\item \emph{Nil and Jacobson radical.} By Theorem~\ref{thm:matrix-nil}, $M_\ell(I_K)$ is nil for every $\ell$. In
  particular $I_K$ is nil, hence Jacobson radical: if $i^N=0$ then $(1-i)^{-1}=\sum_{t<N}i^t$.
\item \emph{Algebraic unitization.} An element $a\cdot1+i$ of $R_K$ satisfies $(t-a)^N$ whenever $i^N=0$.
\end{itemize}
The presentation was fixed before $K$ was chosen, so the same presentation works for every field of positive
characteristic.
\end{proof}

%======================================================================
\section{Unsigned relations in characteristic zero}\label{sec:charzero}

\begin{proof}[Proof of Theorem~C]
Let $A$ be the nonunital algebra over a field $K$ of characteristic zero generated by a finite alphabet $\Sigma$, subject
to finitely many relations $u=v$ and $w=0$ with $u,v,w$ nonempty words. The proof of Lemma~\ref{lem:word-basis} does not
use equal lengths. Relations $u=v$ of different lengths still define a replacement graph, whose classes give a basis.
Hence the nonzero replacement classes of nonempty words form a $K$-basis of $A$.

Let $x=\sum_{a\in\Sigma}[a]$. Expanding,
\[
 x^n=\sum_{w\in\Sigma^n}[w]=\sum_{C}N_n(C)\,[C],
\]
where $C$ runs over the nonzero classes and $N_n(C)$ is the number of words of length $n$ lying in $C$. Each $N_n(C)$ is
a nonnegative integer, and it is positive exactly when $C$ contains a word of length $n$. In characteristic zero, positive
integers are nonzero in $K$. Therefore $x^n=0$ if and only if every word of length $n$ is zero in $A$.

Suppose that $A$ is nil. Then $x^N=0$ for some $N$, so every word of length $N$ vanishes. Every longer word has such a
word as a prefix, and a word with a zero prefix is zero. So every product of $N$ elements of $A$, being a combination of
words of length at least $N$, vanishes. Thus $A^N=0$.

For the presentation of Theorem~A, the words $w_j$ of the proof of Theorem~A are nonzero over every field, by
Lemma~\ref{lem:word-basis}. Their prefixes are nonzero as well, so nonzero words exist in every length. By the argument
above, $x^n\neq0$ for every $n$. In particular the algebra is not nil, and it is infinite-dimensional.
\end{proof}

\begin{remark}
\begin{enumerate}[label=(\arabic*),leftmargin=*]
\item \emph{What Theorem~C does not exclude.} It does not exclude finitely presented nil algebras that are not nilpotent
  in characteristic zero. It shows only that such an algebra cannot be presented by unsigned word relations alone.
  Relations with signs or other coefficients, such as $u+v=0$, binomials $u-\lambda v$, or relations with three or more
  terms, are outside its scope. Whether a finitely presented nil algebra that is not nilpotent exists over $\Q$ is the
  characteristic-zero case of Ufnarovskij's question~\cite[Question~8.1]{Smoktunowicz2008}, and it remains open.
\item \emph{Where the positive-characteristic mechanism lives.} In characteristic $p$, the coefficient of a word class in
  a power of a weighted sum is again a sum of positive integer multiplicities, and Lemma~\ref{lem:symmetric} arranges for
  these to be divisible by $p$. Theorem~C shows that no presentation by unsigned word relations can do this over $\Q$.
\item \emph{Semigroups versus algebras.} An algebra as in Theorem~C is the contracted semigroup algebra of the finitely
  presented semigroup with zero defined by the same relations. Nilness of the semigroup is much weaker than nilness of
  the algebra. Ivanov-Pogodaev and Kanel-Belov construct an infinite finitely presented nilsemigroup satisfying
  $x^9=0$~\cite{IvanovPogodaevKanelBelov2025}. By Theorem~C, its contracted semigroup algebra over any field of
  characteristic zero is not nil.
\item \emph{Priority.} The argument is the elementary observation that the sum of the generators has nonnegative
  integer coefficients on the word basis. We expect the statement to be known in the theory of semigroup
  algebras~\cite{Okninski1991}, and we do not claim it as new. We include it because it identifies exactly why the
  construction of Theorem~A is confined to positive characteristic.
\end{enumerate}
\end{remark}
